\section{What this project is for}

\ProjectName\ is a teaching-oriented implementation of an RV32I pipelined CPU.
Its goal is to help you learn how a modern in-order pipeline is structured and
validated in an HDL, and how the resulting RTL maps onto FPGA-friendly building
blocks (especially synchronous Block RAMs).

\subsection{What you learn by studying or modifying this core}
\begin{itemize}[leftmargin=*]
  \item How a five-stage pipeline is partitioned (Fetch/Decode/Execute/Memory/Writeback) and how stage registers define timing.
  \item How a hazard unit combines forwarding, stalling, and flushing to preserve correctness.
  \item How a Harvard memory map and a custom linker script make bare-metal C programs \emph{actually runnable} on the hardware model.
  \item How to validate ISA behavior with directed tests and self-checking signatures.
  \item HDL coding practice: writing synthesizable RTL (design for synthesis) and writing testbenches for design verification.
  \item Using open-source tools to simulate, verify, and visualize HDL designs (e.g.\ GHDL, GTKWave, DigitalJS).
  \item How to simulate and debug in Xilinx Vivado (XSim), and how to instrument an FPGA build with ILA.
\end{itemize}

\subsection{Textbook lineage}
The RTL follows the structure and naming conventions of the Harris \& Harris
pipelined reference design (stage-suffixed signals such as \texttt{PCF},
\texttt{InstrD}, \texttt{ALUControlE}, \texttt{MemWriteM}, \texttt{ResultSrcW}).
Section~\ref{sec:textbook-deviations} summarizes the intentional deviations made
to support a synthesis-ready, BRAM-friendly implementation and a slightly richer
RV32I subset.

